\begin{table}[tb]
\centering\footnotesize\setlength{\tabcolsep}{3pt}
\caption{予定した全８条件．$+?$は既知小計に使用量不明の生成が残ることを表す．}
\begin{tabular}{lllrrr}
課題 & 条件 & 成果 & 論理token & 生成 & API \\ \hline
写真 & A & 成功 & 135,628 & 16 & 99 \\
写真 & R & 成功 & 135,628 & 16 & 99 \\
精算 & R & 成功 & 211,110 & 24 & 25 \\
精算 & A & 成功 & 211,110 & 24 & 25 \\
返金 & A & 成功 & 102,087 & 13 & 13 \\
返金 & R & 成功 & 102,087 & 13 & 13 \\
書出し & R & 未達 & 433,732 & 22 & 60 \\
書出し & A & 未達 & 320,456$+?$ & 20 & 42 \\
\hline\end{tabular}
\end{table}
成功／評価はAが3/4，Rが3/4であった．課題ごとの成功喪失は0，成功獲得は0である．Aの全仕事tokenは769,281$+?$，Rは882,557であった．
欠測または未評価があるため，固定した全件比較の費用差は未確定とする．固定した開発上の採用基準は満たさなかった．
\begin{table}[tb]
\centering\footnotesize\setlength{\tabcolsep}{3pt}
\caption{全課題の費用寄与．全件既知でないときは全体の寄与を確定しない．}
\begin{tabular}{lrr}
成否の組合せ & 課題数 & R−A token \\ \hline
両側成功 & 3 & 未確定 \\
両側未達 & 1 & 未確定 \\
成否が異なる & 0 & 未確定 \\
片側以上が未評価 & 0 & 未確定 \\
\hline\end{tabular}
\end{table}
CSV課題はRが未達，Aが未達であった．未達の費用も残し，成功部分のみの減少を全体の改善としない．未評価・未実行は未達として補完しない．
